\begin{afterword}
\summaryhead

The post starts from a quarrel: an atheist blames religion for the Inquisition and the Crusades, and a believer answers that atheism is a religion too, or that Stalin was an atheist. Whether Stalin counts as religious, the author shows, depends only on the definition. The real question is why large groups of people have been slaughtered in the name of an idea.

The author's suggested ``ur-mistake'' is the rule of Deuteronomy 13: if anyone, even your brother, tempts you to other gods, do not listen to him; kill him. Stalin and Hitler, the post says, set the same rule. When it feels wrong to argue against any positive claim about the Great Idea, the affective death spiral goes supercritical. This can happen with or without supernatural beliefs, around a political system, a leader or a theory of race, so religion is not the key category. Faith, Sam Harris's target, makes spirals easy to start; they turn deadly when criticism becomes a sin or a crime. The post ends with a rule without exceptions: bad argument gets counterargument, never a bullet.

\responsehead

The post moves a quarrel about a word to a question about history, and its answer is plausible: an idea becomes deadly when criticizing it is forbidden. Its evidence is a passage of scripture and a description of Stalin's and Hitler's rule. The description is accurate in substance. The Soviet state made a hero of a boy said to have denounced his father, and most Gestapo investigations began with denunciations by ordinary Germans.

What the evidence shows is that the rule was present where the killing happened. It does not show that the rule came first and produced the killing, which is what an ``ur-mistake'' would have to be. The post is careful at first: it grants that there may be no single cause and offers its answer as a suggestion. Near the end the hedge is gone, and the spiral simply ``turns much deadlier after criticism becomes a sin, or a gaffe, or a crime.'' No new evidence has been given in between.

Its definition of the supernatural, ``the best distinction I've heard,'' is Richard Carrier's, published earlier the same year. The post does not name him; the author credited Carrier in a later post. The picture of Christianity's defences and of New Agers in spirals ``around stars, trees, magnets, diets, spells, unicorns'' is given without evidence.

The post applies its test to its own side. The first sign of a supercritical spiral on its list is feeling that correcting a flawed argument for evolution would help the creationists. Its closing rule, that no argument for any idea is beyond criticism and that criticism never justifies violence, is stated plainly and without exceptions.

In short: a plausible answer to why ideas become deadly, supported by cases where the suppression of criticism and the killing occurred together, with the step from co-occurrence to cause offered first as a suggestion and later as a plain statement.
\end{afterword}
